\begin{pdfdiagram}[x=1cm,y=1cm]
  % 同心圆只承担区位关系；解释文字移到盘面之外，避免圆弧穿字。
  \draw[BookRule, line width=0.4pt] (5,0) circle (4.5);
  \draw[BookRule, line width=0.4pt] (5,0) circle (3.7);
  \draw[BookRule, line width=0.4pt] (5,0) circle (2.9);
  \draw[BookRule, line width=0.4pt] (5,0) circle (2.1);
  \draw[BookRule, line width=0.4pt] (5,0) circle (1.3);

  \node[hw label, font=\tiny] at (5,4.1) {Zone 0（外）};
  \node[hw label, font=\tiny] at (5,3.3) {Zone 1};
  \node[hw label, font=\tiny] at (5,2.5) {Zone 2};
  \node[hw label, font=\tiny] at (5,1.7) {Zone 3};
  \node[hw label, font=\tiny] at (5,0.9) {Zone 4（内）};

  \node[BookAccent, font=\tiny, anchor=west] at (8.0,3.0) {每道 1200 扇区};
  \node[BookAccent, font=\tiny, anchor=west] at (8.0,1.5) {每道 700 扇区};
  \draw[BookAccent, line width=0.3pt, ->] (7.9,3.0) -- (7.4,3.0) -- (7.4,3.5);
  \draw[BookAccent, line width=0.3pt, ->] (7.9,1.5) -- (6.5,1.5) -- (6.5,1.0);

  \node[hw label, text width=10cm, align=center, anchor=north]
    at (5,-4.85)
    {ZBR：外区每道扇区数多（周长长），内区少。保持全区线密度近似恒定，最大化容量。典型 10--20 个区。};
\end{pdfdiagram}
